% sections/drb.tex -- the dynamic reduced basis, its exactness (Theorem 1) and rebuilds (Proposition 1.1). Uses macros.tex. Executable checks of every statement:
% research/mmhm/reference/check_props.py.
\section{The dynamic reduced basis}\label{sec:drb}

Every cell $x\in K$ carries an \emph{order key} $\ord(x)$; keys of distinct cells
are distinct. For a non-zero chain $c$, $\low(c)$ is the cell of $c$ with the
largest key. Keys do not change between rebuilds (Section~\ref{sec:morse}); an
inserted cell receives a key larger than every key in use. Nothing in this
section depends on how the keys are chosen; \cref{sec:morse} chooses them.

\begin{definition}[DRB]\label{def:drb}
Let $k\ge1$. A \emph{dynamic reduced basis} of $\Bm{k}$ is a finite family of
\emph{slots} $j\in J_k$, each holding a pair $(R_j,V_j)$ with $V_j\in\Cc{k}$
and $R_j\in\Cc{k-1}$, such that
\begin{description}
\item[(I1)] $R_j=\bd_k V_j$ for every $j\in J_k$;
\item[(I2)] the non-zero $R_j$ have pairwise distinct lows;
\item[(I3)] $(V_j)_{j\in J_k}$ is a basis of $\Cc{k}$; in particular $|J_k|=n_k$.
\end{description}
Two index structures are derived from the slots: $\pivot(\ell)=j$ if $R_j\ne0$
and $\low(R_j)=\ell$ (well defined by (I2)), and, for $x\in\Kc{k}$,
$\occ(x)=\{j\in J_k : x\in V_j\}$, the \emph{footprint} of $x$. Both are updated
together with the chains; we do not spell out this bookkeeping.
\end{definition}

Writing $R$ and $V$ for the matrices with columns $R_j$ and $V_j$, (I1) and (I3)
say that $R=\Bm{k}V$ with $V$ invertible, and (I2) says that $R$ is
\emph{reduced}. Unlike the decomposition used for persistence, $V$ is not
required to be triangular with respect to $\ord$.

\paragraph{Insertion.} $\Ins(\sigma)$ for a $k$-simplex $\sigma\notin K$, $k\ge1$,
all of whose facets are in $K$:
\begin{enumerate}
\item give $\sigma$ a key larger than all keys in use; create a slot $j$ with
  $V_j=\{\sigma\}$ and $R_j=\bd\sigma$;
\item while $R_j\ne0$ and $p=\pivot(\low(R_j))$ is defined: $R_j\gets R_j+R_p$,
  $V_j\gets V_j+V_p$;
\item if $R_j\ne0$, set $\pivot(\low(R_j))=j$.
\end{enumerate}
The slots $p_1,\dots,p_m$ added in step~2 form the \emph{reduction chain} of
$\sigma$. We call $\sigma$ \emph{positive} if $R_j=0$ at the end and
\emph{negative} otherwise.

\paragraph{Deletion.} $\Del(\sigma)$ for $\sigma\in\Kc{k}$, $k\ge1$, without
cofaces. Let $O=\occ(\sigma)$ and $Z=\{j\in O: R_j=0\}$.
\begin{description}
\item[Case A ($Z\ne\emptyset$).] Choose any $z\in Z$ (the implementations take
  one with $|V_z|$ minimal, ties broken by the smallest slot index). For every $j\in O\setminus\{z\}$:
  $V_j\gets V_j+V_z$. Remove slot $z$.
\item[Case B ($Z=\emptyset$).] Let $p\in O$ be the slot whose $\low(R_p)$ has
  the smallest key. For every $j\in O\setminus\{p\}$: $R_j\gets R_j+R_p$ and
  $V_j\gets V_j+V_p$. Remove slot $p$ and the entry $\pivot(\low(R_p))$.
\end{description}
Finally remove $\sigma$ from $K$. Inserting or deleting a vertex changes only
$n_0$. \Cref{fig:deletion} carries out both cases on a graph.

\begin{figure}[t]
\centering
% figures/f2_deletion.tex -- the two deletion cases worked on graphs (k = 1).
% Rows of R are vertices with keys 0 < 1 < 2 < 3 (top to bottom, so the low of
% a column is its lowest non-zero entry); rows of V are edges; columns are slots.
% Every entry is computed and checked by reference/figure_examples.py.
\newcommand{\pv}[1]{{\setlength{\fboxsep}{0.6pt}\setlength{\fboxrule}{0.6pt}\fbox{$#1$}}}
\newcommand{\fc}{\cellcolor{black!14}}
\newcommand{\chg}[1]{\boldsymbol{#1}}
\tikzset{
  dvtx/.style={circle, draw=black!85, fill=white, inner sep=0pt, minimum size=3.6mm,
    font=\scriptsize},
  dedg/.style={black!80, semithick},
  dsig/.style={black, very thick, dash pattern=on 2.4pt off 1.5pt},
}
\newcommand{\whynote}[1]{{\scriptsize\raggedright #1\par}}
{\footnotesize\renewcommand{\arraystretch}{1.06}%
%
% =========================== (a) Case A ===================================
\begin{minipage}[t]{0.485\textwidth}
\textbf{(a) Case A:} $\sigma=01$ lies on a cycle\par\smallskip
\begin{tabular}{@{}c@{\hspace{0.7em}}l@{\hspace{1.1em}}l@{}}
\begin{tikzpicture}[baseline=(current bounding box.center)]
  \node[dvtx] (a0) at (0,0) {0}; \node[dvtx] (a1) at (1.1,0) {1};
  \node[dvtx] (a2) at (0.55,0.95) {2};
  \draw[dedg] (a0) -- (a2) (a1) -- (a2);
  \draw[dsig] (a0) -- node[below=-1pt, font=\scriptsize] {$\sigma$} (a1);
  \node[font=\scriptsize] at (0.55,-0.62) {$\beta=(1,1)$};
\end{tikzpicture}
& $V=\begin{array}{@{}r|ccc@{}}
    & s_1 & s_2 & s_3\\ \hline
    \fc 01 & \fc 1 & \fc 0 & \fc 1\\
    02 & 0 & 1 & 1\\
    12 & 0 & 0 & 1\end{array}$
& $R=\begin{array}{@{}r|ccc@{}}
    & s_1 & s_2 & s_3\\ \hline
    0 & 1 & 1 & 0\\
    1 & \pv{1} & 0 & 0\\
    2 & 0 & \pv{1} & 0\end{array}$\\[16pt]
$\Big\downarrow$ & \multicolumn{2}{@{}l@{}}{\parbox[c]{5.75cm}{\scriptsize\raggedright
    $O=\occ(\sigma)=\{s_1,s_3\}$ (the $1$s in the shaded row); $R_{s_3}=0$, so
    $Z=\{s_3\}$ and $z=s_3$.\\
    $V_{s_1}\gets V_{s_1}+V_{s_3}$; remove slot $s_3$ and row $\sigma$.}}\\[14pt]
\begin{tikzpicture}[baseline=(current bounding box.center)]
  \node[dvtx] (b0) at (0,0) {0}; \node[dvtx] (b1) at (1.1,0) {1};
  \node[dvtx] (b2) at (0.55,0.95) {2};
  \draw[dedg] (b0) -- (b2) (b1) -- (b2);
  \node[font=\scriptsize] at (0.55,-0.62) {$\beta=(1,0)$};
\end{tikzpicture}
& $V=\begin{array}{@{}r|cc@{}}
    & s_1 & s_2\\ \hline
    02 & \chg{1} & 1\\
    12 & \chg{1} & 0\end{array}$
& $R=\begin{array}{@{}r|cc@{}}
    & s_1 & s_2\\ \hline
    0 & 1 & 1\\
    1 & \pv{1} & 0\\
    2 & 0 & \pv{1}\end{array}$
\end{tabular}\par\medskip
\whynote{\textbf{Why a cycle.} $\bd V_{s_3}=R_{s_3}=0$, so the addition removes
  $\sigma$ from $V_{s_1}$ and leaves every $R_j$, hence every low, as it was.
  $n_1$ drops and $r_1=2$ stays: $\beta_1$ falls by one.}
\end{minipage}\hfill
%
% =========================== (b) Case B ===================================
\begin{minipage}[t]{0.485\textwidth}
\textbf{(b) Case B:} $\sigma=03$ lies on no cycle\par\smallskip
\begin{tabular}{@{}c@{\hspace{0.7em}}l@{\hspace{1.1em}}l@{}}
\begin{tikzpicture}[baseline=(current bounding box.center)]
  \node[dvtx] (c3) at (0.6,1.0) {3};
  \node[dvtx] (c0) at (0,0) {0}; \node[dvtx] (c1) at (0.6,0) {1};
  \node[dvtx] (c2) at (1.2,0) {2};
  \draw[dedg] (c1) -- (c3) (c2) -- (c3);
  \draw[dsig] (c0) -- node[left=-1pt, font=\scriptsize] {$\sigma$} (c3);
  \node[font=\scriptsize] at (0.6,-0.62) {$\beta=(1,0)$};
\end{tikzpicture}
& $V=\begin{array}{@{}r|ccc@{}}
    & s_1 & s_2 & s_3\\ \hline
    \fc 03 & \fc 1 & \fc 1 & \fc 1\\
    13 & 0 & 1 & 0\\
    23 & 0 & 0 & 1\end{array}$
& $R=\begin{array}{@{}r|ccc@{}}
    & s_1 & s_2 & s_3\\ \hline
    0 & 1 & 1 & 1\\
    1 & 0 & \pv{1} & 0\\
    2 & 0 & 0 & \pv{1}\\
    3 & \pv{1} & 0 & 0\end{array}$\\[16pt]
$\Big\downarrow$ & \multicolumn{2}{@{}l@{}}{\parbox[c]{5.75cm}{\scriptsize\raggedright
    $O=\{s_1,s_2,s_3\}$ and every $R_j\ne0$, so $Z=\emptyset$; the lows are
    $3,1,2$, and $p=s_2$ has the smallest.\\
    $(R_j,V_j)\gets(R_j+R_{s_2},V_j+V_{s_2})$ for $j=s_1,s_3$; remove slot $s_2$
    and row $\sigma$.}}\\[14pt]
\begin{tikzpicture}[baseline=(current bounding box.center)]
  \node[dvtx] (d3) at (0.6,1.0) {3};
  \node[dvtx] (d0) at (0,0) {0}; \node[dvtx] (d1) at (0.6,0) {1};
  \node[dvtx] (d2) at (1.2,0) {2};
  \draw[dedg] (d1) -- (d3) (d2) -- (d3);
  \node[font=\scriptsize] at (0.6,-0.62) {$\beta=(2,0)$};
\end{tikzpicture}
& $V=\begin{array}{@{}r|cc@{}}
    & s_1 & s_3\\ \hline
    13 & \chg{1} & \chg{1}\\
    23 & 0 & 1\end{array}$
& $R=\begin{array}{@{}r|cc@{}}
    & s_1 & s_3\\ \hline
    0 & \chg{0} & \chg{0}\\
    1 & \chg{1} & \chg{1}\\
    2 & 0 & \pv{1}\\
    3 & \pv{1} & 0\end{array}$
\end{tabular}\par\medskip
\whynote{\textbf{Why the smallest low.} Adding $R_{s_1}=\{0,3\}$, whose low $3$
  is the largest, would instead give $\{1,3\}$ and $\{2,3\}$: two lows $3$,
  violating (I2). $R_{s_2}$ lies entirely above the other lows, so adding it
  moves none. $n_1$ and $r_1$ both drop: $\beta_0$ rises by one.}
\end{minipage}}%

\caption{The two cases of $\Del(\sigma)$, worked on graphs ($k=1$). Columns are
slots. The rows of $V$ are edges; the rows of $R$ are vertices in key order from
top to bottom, so the low of a column is its lowest non-zero entry (boxed). The
shaded row of $V$ is $\sigma$, and its $1$s mark the footprint $O$. Bold entries
are those changed by the update; $\beta=(\beta_0,\beta_1)$.
(a)~A footprint slot is a cycle ($R_{s_3}=0$); adding it to the other footprint
slot clears $\sigma$ from $V$ and leaves $R$ unchanged.
(b)~No footprint slot is a cycle; the slot with the smallest low is added to the
others and removed, and no low moves.
Every entry is recomputed and checked by \texttt{reference/figure\_examples.py}.}
\label{fig:deletion}
\end{figure}

\begin{theorem}[Exactness]\label{thm:exact}
Let $k\ge1$ and let $(R_j,V_j)_{j\in J_k}$ satisfy (I1)--(I3) for $\Bm{k}(K)$,
for some injective key function.
\begin{enumerate}
\item[(a)] $r_k=\#\{j\in J_k : R_j\ne0\}$, and $\{V_j : R_j=0\}$ is a basis of
  $\Zc{k}(K)$.
\item[(b)] Let $\sigma\notin K$ be a $k$-simplex whose facets are all in $K$.
  The loop of $\Ins(\sigma)$
  terminates, and afterwards the slots satisfy (I1)--(I3) for
  $\Bm{k}(K\cup\{\sigma\})$; every family satisfying (I1)--(I3) for
  $\Bm{k+1}(K)$ satisfies them, unchanged, for $\Bm{k+1}(K\cup\{\sigma\})$.
\item[(c)] Let $\sigma\in\Kc{k}$ have no coface. Then $O\ne\emptyset$; after
  $\Del(\sigma)$, in Case A and in Case B, the slots satisfy (I1)--(I3) for
  $\Bm{k}(K\setminus\{\sigma\})$; every family satisfying (I1)--(I3) for
  $\Bm{k+1}(K)$ satisfies them, unchanged, for $\Bm{k+1}(K\setminus\{\sigma\})$.
\item[(d)] Inserting a positive $k$-simplex increases $\beta_k$ by one; inserting
  a negative one decreases $\beta_{k-1}$ by one. A deletion in Case A decreases
  $\beta_k$ by one; a deletion in Case B increases $\beta_{k-1}$ by one. No
  other Betti number changes. Moreover, $\Del(\sigma)$ is in Case A if and only
  if $\sigma$ lies in the support of a $k$-cycle of $K$, and $\Ins(\sigma)$ is
  positive if and only if $\sigma$ lies in the support of a $k$-cycle of
  $K\cup\{\sigma\}$.
\end{enumerate}
\end{theorem}

\begin{proof}
(a) By (I1) and (I3), $R=\Bm{k}V$ with $V$ invertible, so $\rank R=r_k$. The
non-zero columns of $R$ are linearly independent: in a non-empty sum of them,
let $R_i$ be the summand whose low $\ell$ has the largest key; every other
summand has its low, hence all its cells, strictly below $\ell$, so $\ell$
survives in the sum. Hence $r_k$ is the number of non-zero $R_j$. Each $V_j$ with
$R_j=0$ is a cycle; these chains are independent as part of a basis, and there
are $n_k-r_k=\dim\Zc{k}$ of them.

(b) Let $K'=K\cup\{\sigma\}$. Throughout the loop, $R_j=\bd V_j$ (true at the
start and preserved by linearity because $R_p=\bd V_p$) and $V_j=\sigma+w$ with
$w\in\Cc{k}(K)$. An addition is performed only when $\low(R_p)=\low(R_j)=\ell$;
all cells of both chains have keys at most $\ord(\ell)$ and $\ell$ cancels, so
afterwards $R_j=0$ or $\low(R_j)$ is strictly smaller. Hence the loop
terminates, after at most one addition per non-zero slot. On exit either
$R_j=0$ or $\low(R_j)$ is no pivot, so recording the pivot preserves (I2); (I1)
holds. Since $\Cc{k}(K')=\Cc{k}(K)\oplus\Z\sigma$ and the old $V_i$ form a basis
of $\Cc{k}(K)$, the old chains together with $V_j=\sigma+w$ form a basis of
$\Cc{k}(K')$: they are $n_k+1$ chains spanning $\sigma=V_j+w$. For $\Bm{k+1}$:
the chain group $\Cc{k+1}$ is unchanged, each $R_i=\bd V_i$ is supported on facets
of $(k+1)$-simplices of $K'$, none of which is $\sigma$
(Observation~\ref{obs:structural}), and keys of existing cells are unchanged;
so (I1)--(I3) continue to hold.

(c) Since the $V_j$ span $\Cc{k}$, $\sigma$ is a sum of some of them, so some
$V_j$ contains $\sigma$: $O\ne\emptyset$. Each operation $V_j\gets V_j+V_q$ with
$j\ne q$ is an invertible change of basis, so throughout the loop the family
remains a basis of $\Cc{k}(K)$. For $j\in O\setminus\{q\}$ (with $q=z$ in Case A,
$q=p$ in Case B), both $V_j$ and $V_q$ contain $\sigma$, so $V_j+V_q$ does not;
slots outside $O$ never contained $\sigma$. After the loop, $V_q$ is therefore
the only chain containing $\sigma$. Removing slot $q$ leaves $n_k-1$ independent
chains, all in $\Cc{k}(K\setminus\{\sigma\})$, which has dimension $n_k-1$:
(I3) holds.
\emph{Case A.} Because $\bd V_z=R_z=0$, the additions leave every $R_j$
unchanged, so (I1) and (I2) hold; the removed slot had $R_z=0$ and owned no
pivot.
\emph{Case B.} Every $R_j$ with $j\in O$ is non-zero, and by (I2) the lows are
distinct, so $p$ is well defined and $\low(R_p)$ is strictly smaller than
$\low(R_j)$ for $j\in O\setminus\{p\}$. All cells of $R_p$ have keys at most
$\ord(\low(R_p))<\ord(\low(R_j))$, so $\low(R_j)\notin R_p$ and
$\low(R_j+R_p)=\low(R_j)$; in particular $R_j+R_p\ne0$. Lows are unchanged and
the low of $R_p$ is released with slot $p$, so (I2) holds; (I1) holds by
linearity. (The minimal choice of $p$ is what guarantees this: adding an $R_p$
with a larger low would move lows and require further reduction.)
For $\Bm{k+1}$: $\sigma$ has no coface, so no $R_i=\bd V_i$ of dimension $k+1$
contains $\sigma$ and $\sigma$ is not a pivot there; deleting the zero row
$\sigma$ leaves (I1)--(I3) intact.

(d) By (a), $r_k$ increases by one exactly when the inserted $\sigma$ is negative,
and decreases by one exactly in Case B; $n_k$ changes by $\pm1$; $r_{k+1}$ is
unchanged because only a zero row of $\Bm{k+1}$ is added or removed
(Observation~\ref{obs:structural}); all other $n_i$ and $r_i$ are unchanged.
Formula~\eqref{eq:betti} gives the stated changes. If
$z\in Z$, then $V_z$ is a cycle containing $\sigma$. Conversely, if a cycle
$c\in\Zc{k}(K)$ contains $\sigma$, write $c$ in the basis
$\{V_j:R_j=0\}$ of $\Zc{k}(K)$ from (a); some $V_j$ in this expansion contains
$\sigma$, and that $j$ lies in $Z$. For insertion, if $\sigma$ is positive then
the final $V_j$ is a cycle of $K'$ containing $\sigma$. Conversely, if a cycle
$c$ of $K'$ contains $\sigma$, then $\bd\sigma=\bd(c+\sigma)$ lies in the
column space of $\Bm{k}(K)$, so $r_k$ does not increase, and by (a) the new slot
has $R_j=0$.
\end{proof}

Vertices need no slots: inserting a vertex increases $n_0$ and $\beta_0$ by
one, and a deletable vertex is isolated, so deleting it decreases both by one
(Observation~\ref{obs:structural} and~\eqref{eq:betti}). A batch is processed
one atomic edit at a time, so Theorem~\ref{thm:exact} makes the Betti numbers
exact after every update. The choice of $z$ in Case A
is free; it only affects sparsity. Theorem~\ref{thm:exact} holds for every key
function, so the order chosen in Section~\ref{sec:morse} is a matter of cost
alone.

\paragraph{Initialisation and rebuilds.}
The DRB is initialised, and periodically rebuilt, by the static reduction with
clearing~\cite{chen2011twist}, tracking $V$: for $k=d,d-1,\dots,1$, process the
columns of $\Bm{k}$ in increasing key order; if the $k$-simplex $\tau$ is the low
of a non-zero reduced column $R$ of $\Bm{k+1}$, \emph{clear} it ($R_\tau:=0$,
$V_\tau:=R$); otherwise reduce it as in steps~2--3 of $\Ins$, starting from
$V_\tau=\{\tau\}$, $R_\tau=\bd\tau$.

\begin{subproposition}[Rebuild]\label{cor:rebuild}
For every injective key function, the static reduction with clearing produces
slots satisfying (I1)--(I3) in every dimension. Moreover, every $V_\tau$ consists
of $\tau$ and cells with smaller keys. The same holds for the reduction without
clearing. In particular the ranks it reports are exact even when $V$ is not
stored.
\end{subproposition}

\begin{proof}
We show $\low(V_\tau)=\tau$ for every column $\tau$, by induction along the
processing order. A non-cleared $\tau$ ends with
$V_\tau=\{\tau\}+\sum V_p$ over pivot owners $p$ processed earlier, hence with
smaller keys, and by induction every cell of $V_p$ has key at most $\ord(p)$. A
cleared $\tau$ has $V_\tau=R$ with $\low(R)=\tau$. Ordering rows and columns by
key, $V$ is therefore upper unitriangular, hence invertible: (I3). (I1) holds
for non-cleared columns by linearity, and for cleared ones because $R=\bd V'$
for the slot $V'$ of $\Bm{k+1}$ that produced it (induction over dimensions), so
$\bd V_\tau=\bd\bd V'=0=R_\tau$. (I2): a non-cleared column ends with $R_\tau=0$ or
with a low that no earlier column owns; cleared columns are zero. Without
clearing every column is of the non-cleared kind and the same argument applies.
The last statement follows from Theorem~\ref{thm:exact}(a).
\end{proof}

Note that the correctness of clearing here follows from the invariants alone;
we do not need the classical fact that a cleared column would have reduced to
zero. After deletions, however, $V$ is no longer triangular (Case A adds
$V_z$, whose low may exceed that of $V_j$), so the DRB maintains ranks and a cycle
basis but not a persistence pairing, and we claim none. This is the difference
from vineyard updates~\cite{cohensteiner2006vines}, which keep $V$ triangular
for a changing order.

\paragraph{Rebuild policy.} \method{mmhm} rebuilds with the Morse order of Section~\ref{sec:morse}; it
rebuilds after update $t$ if $t-t_{\mathrm{last}}\ge m$ or
$\nnz(R)+\nnz(V)>(1+\alpha)\max(\nnz_{\mathrm{last}},|K|)$, where
$\nnz_{\mathrm{last}}$ is the value right after the previous rebuild (defaults
$m=\infty$, $\alpha=1$), and the rebuild is charged to update $t$.
\method{drb-id} uses insertion keys and never rebuilds; \method{mmhm-norb} uses
the Morse order at the initial build only; \method{drb-id-rb} applies the same
policy as \method{mmhm} but rebuilds in insertion order
(Section~\ref{sec:exp-methods}).
